\documentclass[11pt,a4paper]{whatsnote}
\usepackage{siunitx}
\coverset{title=Experimental Notes,author=Your Name,afill=Science notebook,date=\today,extinfo=Thermal cycling}
\begin{document}
\maketitle
\section{Question}
Does repeated cycling between \SI{-20}{\celsius} and \SI{80}{\celsius} change
the sample's storage modulus? Define the acceptance criterion before looking
at the final measurements.
\section{Setup}
Instrument: DMA-000. Sample: polymer strip, batch 00. Record calibration file,
fixture, atmosphere, ramp rate and operator here. Keep a separate copy of raw
instrument exports; this note records decisions and observations.
\section{Run log}
\begin{description}
\item[08:40] Calibration check passed; baseline displacement was stable.
\item[09:15] First cycle completed. A small edge curl appeared after the hot hold.
\item[10:30] Second cycle paused at the cold plateau to check the thermocouple.
\end{description}
\section{Observations and decisions}
Write what changed, why the change was made, and which raw file supports it.
\subsection*{Next action}
Repeat the cycle with a fresh specimen and compare the fitted modulus at the
same frequency.\\[2em]
\noindent\rule{\linewidth}{0.4pt}
\end{document}
